\section{Introduction}
\label{sec:introduction}

Construction remains one of the most hazardous industries worldwide. The spatially dynamic nature of construction sites creates persistent risks because workers, heavy equipment, materials, temporary structures, and active work zones continuously change their positions and interactions. Effective safety monitoring must therefore extend beyond recognizing objects or unsafe behavior in isolated video frames and should support continuous interpretation of where workers and equipment are located, how they move, and whether their spatial relationships indicate hazardous proximity.

Positioning technologies such as GPS/GNSS, RFID, UWB, and BLE have been used to localize construction resources and maintain persistent identities. These technologies can provide direct positional measurements, but they require tags or devices to be attached to each tracked entity and introduce operational burdens associated with deployment, calibration, power, communication coverage, maintenance, and retrieval. Positional sensors also do not inherently provide visual information about personal protective equipment, surrounding work context, or observable worker behavior. For these reasons, considerable research has investigated non-contact construction monitoring using existing CCTV infrastructure and computer vision.

Recent advances in deep learning have enabled automated hardhat and harness detection, unsafe-behavior recognition, and worker and equipment tracking from construction videos. Pose estimation and person re-identification have also been introduced to reduce track fragmentation and identity switches when workers wear visually similar personal protective equipment. Nevertheless, image coordinates are not physical site coordinates. Two-dimensional observations from different cameras are governed by different projective geometries and cannot be compared directly. The same pixel separation can correspond to substantially different physical distances depending on object depth and camera viewpoint. Single-camera trajectories are further interrupted by occlusion, missed detections, illumination changes, and limited field of view.

Multi-camera methods address some of these limitations by connecting identities across views. Construction-oriented studies have used re-identification for safety-compliance monitoring, matched earthmoving equipment across non-overlapping cameras, and combined appearance, location, and motion cues for cross-camera worker association. BIM-integrated monitoring has also provided spatial context by visualizing camera-derived observations in three-dimensional building models. However, these approaches generally treat cross-camera identity matching and physical localization as separate problems. They do not consistently represent each camera observation as a metric 3D estimate with explicit uncertainty in a common construction-site coordinate system.

Three limitations therefore remain. First, most construction multi-camera tracking methods preserve identity continuity without transforming all observations into a common metric 3D site frame. Consequently, trajectories and inter-object distances are not consistently interpretable in physical units across cameras. Second, geometry-based approaches commonly rely on planar homographies, simplified ground-plane assumptions, or scene-specific Euclidean association thresholds. Such assumptions are restrictive in irregular and evolving construction environments, and they do not account explicitly for the fact that localization uncertainty varies with camera pose, object distance, view direction, occlusion, and depth estimation quality. Third, previous studies often evaluate detection, tracking, localization, or visualization separately. End-to-end evidence linking camera-to-map registration, metric 3D lifting, uncertainty-aware global association, field deployment, and safety-relevant spatial analysis remains limited.

To address these limitations, this study proposes SAFE-3D, standing for Site-Aligned Association and Fusion with Uncertainty Estimation in 3D. SAFE-3D uses a UAV-reconstructed site map as a common metric world frame for multiple fixed CCTV cameras. Each CCTV camera is registered to the site map, and each detection is lifted into metric 3D using a representative ground-contact anchor and scale-corrected depth. An anisotropic covariance is assigned to each observation to represent depth- and viewpoint-dependent uncertainty. The same uncertainty representation is then used for inverse-covariance fusion, Mahalanobis-distance gating, Kalman filtering, and motion-gated tracklet re-linking. The objective is not to introduce a new detector or single-camera tracker, but to transform existing camera-specific outputs into physically interpretable global trajectories in a shared construction-site coordinate system.

The contributions of this study correspond directly to the three limitations identified above:

\begin{enumerate}[leftmargin=*]
    \item \textbf{UAV-referenced metric 3D localization across multiple CCTV cameras.} A UAV-reconstructed site map establishes a common metric frame, and each fixed CCTV camera is registered to this map. Camera-specific detections are converted into site-coordinate 3D observations without requiring a separately hand-constructed ground-plane homography for every camera.

    \item \textbf{Uncertainty-aware multi-view fusion and global association.} Depth- and viewpoint-dependent anisotropic covariance is modeled during 2D-to-3D lifting. The covariance is used consistently for observation fusion, statistical gating, trajectory estimation, and tracklet re-linking, avoiding scene-specific Euclidean distance threshold tuning.

    \item \textbf{End-to-end validation from controlled experiments to operating construction sites.} Detection, camera registration, metric localization, component ablation, cross-camera association, trajectory reconstruction, and proximity analysis are evaluated across four complementary datasets. The analysis further quantifies the camera dependence and depth ambiguity of image-plane proximity cues.
\end{enumerate}

The remainder of this paper is organized as follows. Section~\ref{sec:related_work} reviews vision-based construction monitoring, sensor-based localization, multi-camera tracking, and spatial registration. Section~\ref{sec:method} presents SAFE-3D. Section~\ref{sec:results} reports the experimental results. Section~\ref{sec:discussion} discusses the findings, practical implications, and limitations. Section~\ref{sec:conclusion} concludes the paper.
